% ═══════════════════════════════════════════════════════════════
% LERNHEFT: Bewegungslehre - Kinematik Klasse 10E Gymnasium
% Klausurvorbereitung - Vollständige Übersicht
% ═══════════════════════════════════════════════════════════════

\documentclass[11pt, a4paper]{article}

% ═══════════════════════════════════════════════════════════════
% PAKETE
% ═══════════════════════════════════════════════════════════════

\usepackage[utf8]{inputenc}
\usepackage[T1]{fontenc}
\usepackage[ngerman]{babel}
\usepackage{helvet}
\renewcommand{\familydefault}{\sfdefault}

\usepackage[margin=2cm, top=2.5cm, bottom=2cm]{geometry}
\usepackage{amsmath, amssymb}
\usepackage{siunitx}
\sisetup{locale = DE, per-mode = symbol}

\usepackage{graphicx}
\usepackage{tikz}
\usetikzlibrary{arrows.meta, calc, patterns}
\usepackage{pgfplots}
\pgfplotsset{compat=1.18}

\usepackage{enumitem}
\usepackage{tabularx}
\usepackage{booktabs}
\usepackage{array}
\usepackage{multicol}

\usepackage{xcolor}
\usepackage{tcolorbox}
\tcbuselibrary{skins, breakable}
\usepackage{fancyhdr}
\usepackage{lastpage}
\usepackage{hyperref}

% ═══════════════════════════════════════════════════════════════
% FARBEN
% ═══════════════════════════════════════════════════════════════

\definecolor{physik}{HTML}{2c5aa0}
\definecolor{success}{HTML}{2a7a4b}
\definecolor{warning}{HTML}{b35c00}
\definecolor{grau}{HTML}{888888}
\definecolor{hellgrau}{HTML}{f5f5f5}
\definecolor{hellblau}{HTML}{e8f0f8}

% ═══════════════════════════════════════════════════════════════
% BOXEN
% ═══════════════════════════════════════════════════════════════

\newtcolorbox{formelbox}[1][]{
    colback=hellblau,
    colframe=physik,
    boxrule=1pt,
    arc=2pt,
    left=8pt, right=8pt, top=8pt, bottom=8pt,
    fonttitle=\bfseries,
    #1
}

\newtcolorbox{merksatz}[1][]{
    colback=white,
    colframe=success,
    boxrule=1.5pt,
    arc=2pt,
    left=8pt, right=8pt, top=6pt, bottom=6pt,
    fonttitle=\bfseries,
    title=Merke,
    #1
}

\newtcolorbox{warnbox}[1][]{
    colback=white,
    colframe=warning,
    boxrule=1.5pt,
    arc=2pt,
    left=8pt, right=8pt, top=6pt, bottom=6pt,
    fonttitle=\bfseries,
    title=Achtung,
    #1
}

\newtcolorbox{beispielbox}[1][]{
    colback=hellgrau,
    colframe=grau,
    boxrule=0.5pt,
    arc=2pt,
    left=8pt, right=8pt, top=6pt, bottom=6pt,
    fonttitle=\bfseries,
    title=Beispiel,
    #1
}

% ═══════════════════════════════════════════════════════════════
% KOPF- UND FUSSZEILE
% ═══════════════════════════════════════════════════════════════

\pagestyle{fancy}
\fancyhf{}
\renewcommand{\headrulewidth}{0.4pt}
\renewcommand{\footrulewidth}{0pt}

\lhead{\small\textcolor{physik}{\textbf{Physik 10E}} -- Kinematik}
\rhead{\small Klausurvorbereitung}
\cfoot{\small\thepage/\pageref{LastPage}}

% ═══════════════════════════════════════════════════════════════
% BEFEHLE
% ═══════════════════════════════════════════════════════════════

\newcommand{\wichtig}[1]{\textcolor{physik}{\textbf{#1}}}
\newcommand{\einheit}[1]{\,\si{#1}}
\newcommand{\luecke}[1]{\underline{\hspace{#1}}}

% ═══════════════════════════════════════════════════════════════
% DOKUMENT
% ═══════════════════════════════════════════════════════════════

\begin{document}

% ═══════════════════════════════════════════════════════════════
% KOPFBEREICH (ohne Titelseite)
% ═══════════════════════════════════════════════════════════════

\begin{center}
    {\Large\textbf{\textcolor{physik}{Lernheft Bewegungslehre -- Kinematik}}}
\end{center}

\vspace{0.5em}

\begin{tcolorbox}[colback=white, colframe=red!70!black, boxrule=1.5pt, arc=2pt, left=8pt, right=8pt, top=6pt, bottom=6pt]
\textcolor{red!70!black}{\textbf{Hinweis:}} Dieses Lernheft dient der \textbf{Wiederholung} und als \textbf{Grundlage für das erweiterte Selbststudium}. Es erhebt \textbf{keinen Anspruch auf Vollständigkeit} bezüglich der möglichen Klausuraufgabenstellungen.
\end{tcolorbox}

\vspace{0.5em}

\tableofcontents

\vspace{1em}

% ═══════════════════════════════════════════════════════════════
% KAPITEL 1: GLEICHFÖRMIGE BEWEGUNG
% ═══════════════════════════════════════════════════════════════

\section{Gleichförmige Bewegung}

\subsection{Definition und Grundformel}

\begin{merksatz}
Eine \wichtig{gleichförmige Bewegung} liegt vor, wenn ein Körper in gleichen Zeitabschnitten gleiche Wegstrecken zurücklegt. Die \wichtig{Geschwindigkeit $v$} ist dabei konstant.
\end{merksatz}

\begin{formelbox}[title=Grundformel der gleichförmigen Bewegung]
\begin{center}
{\Large $v = \dfrac{s}{t}$ \hspace{2cm} $s = v \cdot t$ \hspace{2cm} $t = \dfrac{s}{v}$}
\end{center}
\vspace{0.3em}
\begin{tabular}{lll}
$v$ & = & Geschwindigkeit in $\si{\metre\per\second}$ oder $\si{\kilo\metre\per\hour}$ \\
$s$ & = & Weg (Strecke) in $\si{\metre}$ oder $\si{\kilo\metre}$ \\
$t$ & = & Zeit in $\si{\second}$ oder $\si{\hour}$ \\
\end{tabular}
\end{formelbox}

\subsection{Bezugssystem}

\begin{merksatz}
Ein \wichtig{Bezugssystem} ist der Standpunkt, von dem aus eine Bewegung beschrieben wird. Die Geschwindigkeit eines Körpers hängt vom gewählten Bezugssystem ab!
\end{merksatz}

\begin{beispielbox}
\textbf{Beispiel: Passagier im Zug}

Ein Zug fährt mit $v_{\text{Zug}} = \SI{100}{\kilo\metre\per\hour}$ nach Osten. Ein Passagier geht mit $v_{\text{Passagier}} = \SI{5}{\kilo\metre\per\hour}$ durch den Zug.

\begin{itemize}[leftmargin=1em]
    \item \textbf{Bezugssystem Zug:} Der Passagier bewegt sich mit $v = \SI{5}{\kilo\metre\per\hour}$
    \item \textbf{Bezugssystem Erde (Passagier geht nach vorne):} $v = \SI{100}{\kilo\metre\per\hour} + \SI{5}{\kilo\metre\per\hour} = \SI{105}{\kilo\metre\per\hour}$
    \item \textbf{Bezugssystem Erde (Passagier geht nach hinten):} $v = \SI{100}{\kilo\metre\per\hour} - \SI{5}{\kilo\metre\per\hour} = \SI{95}{\kilo\metre\per\hour}$
\end{itemize}
\end{beispielbox}

\begin{warnbox}
Geschwindigkeitsangaben sind nur sinnvoll, wenn das Bezugssystem klar ist. ``Der Passagier bewegt sich'' -- relativ zu was?
\end{warnbox}

\subsection{Einheitenumrechnung}

\begin{formelbox}[title=Umrechnung km/h $\leftrightarrow$ m/s]
\begin{center}
{\large $\si{\kilo\metre\per\hour} \xrightarrow{\div 3{,}6} \si{\metre\per\second}$ \hspace{3cm} $\si{\metre\per\second} \xrightarrow{\times 3{,}6} \si{\kilo\metre\per\hour}$}
\end{center}
\end{formelbox}

\begin{beispielbox}
\textbf{Kopfrechenbare Umrechnungen:}

\begin{center}
\begin{tabular}{|c|c|c|c|c|c|c|}
\hline
\textbf{km/h} & 36 & 72 & 90 & 108 & 144 & 180 \\
\hline
\textbf{m/s} & 10 & 20 & 25 & 30 & 40 & 50 \\
\hline
\end{tabular}
\end{center}
\end{beispielbox}

\subsection{Weg-Zeit-Diagramm ($s$-$t$-Diagramm)}

\begin{minipage}{0.55\textwidth}
\begin{tikzpicture}[scale=0.9]
    \begin{axis}[
        xlabel={$t$ in s},
        ylabel={$s$ in m},
        xmin=0, xmax=6,
        ymin=0, ymax=50,
        grid=major,
        width=8cm, height=6cm,
        axis lines=left,
    ]
    \addplot[very thick, physik, domain=0:5] {10*x};
    \addplot[very thick, success, domain=0:5] {5*x};
    \addplot[very thick, warning, domain=0:5] {20};
    \end{axis}
\end{tikzpicture}
\end{minipage}
\hfill
\begin{minipage}{0.4\textwidth}
\textbf{Interpretation:}
\begin{itemize}[leftmargin=1em]
    \item \textcolor{physik}{\textbf{Blau:}} $v = \SI{10}{\metre\per\second}$
    \item \textcolor{success}{\textbf{Grün:}} $v = \SI{5}{\metre\per\second}$
    \item \textcolor{warning}{\textbf{Orange:}} Stillstand ($v = 0$)
\end{itemize}

\vspace{0.5em}
\textbf{Steigung = Geschwindigkeit}
$$v = \frac{\Delta s}{\Delta t}$$
\end{minipage}

\subsection{Geschwindigkeit-Zeit-Diagramm ($v$-$t$-Diagramm)}

\begin{minipage}{0.55\textwidth}
\begin{tikzpicture}[scale=0.9]
    \begin{axis}[
        xlabel={$t$ in s},
        ylabel={$v$ in m/s},
        xmin=0, xmax=6,
        ymin=0, ymax=15,
        grid=major,
        width=8cm, height=5cm,
        axis lines=left,
    ]
    \addplot[very thick, physik, domain=0:5] {10};
    \addplot[very thick, success, domain=0:5] {5};
    \fill[physik, opacity=0.2] (0,0) rectangle (5,10);
    \end{axis}
\end{tikzpicture}
\end{minipage}
\hfill
\begin{minipage}{0.4\textwidth}
\textbf{Interpretation:}
\begin{itemize}[leftmargin=1em]
    \item Horizontale Linie = konstante Geschwindigkeit
    \item \wichtig{Fläche unter dem Graphen = zurückgelegter Weg}
\end{itemize}

$$s = v \cdot t = \text{Fläche}$$
\end{minipage}

\subsection{Die Fläche unter dem $v$-$t$-Graphen}

\begin{merksatz}
Die \wichtig{Fläche unter dem $v$-$t$-Graphen} entspricht immer dem \wichtig{zurückgelegten Weg} -- unabhängig von der Bewegungsart!
\end{merksatz}

\textbf{Begründung:} Der Weg ergibt sich aus Geschwindigkeit $\times$ Zeit. Dies entspricht genau der Fläche eines Rechtecks (oder eines Dreiecks bei Beschleunigung).

\begin{beispielbox}[title=Beispiel: Weg aus $v$-$t$-Diagramm berechnen]
\textbf{Situation:} Ein Auto beschleunigt $\SI{4}{\second}$ lang mit $a = \SI{5}{\metre\per\second\squared}$ und fährt dann $\SI{4}{\second}$ gleichförmig weiter.

\begin{center}
\begin{tikzpicture}[scale=0.8]
    \begin{axis}[
        xlabel={$t$ in s},
        ylabel={$v$ in m/s},
        xmin=0, xmax=10,
        ymin=0, ymax=25,
        grid=major,
        width=10cm, height=5cm,
        axis lines=left,
    ]
    % Phase 1: Beschleunigung (0-4s)
    \addplot[very thick, physik, domain=0:4] {5*x};
    % Phase 2: Gleichförmig (4-8s)
    \addplot[very thick, success, domain=4:8] {20};
    % Flächen markieren
    \fill[physik, opacity=0.2] (0,0) -- (4,0) -- (4,20) -- cycle;
    \fill[success, opacity=0.2] (4,0) rectangle (8,20);
    % Beschriftungen
    \node at (2, 5) {\small $A_1$};
    \node at (6, 10) {\small $A_2$};
    \end{axis}
\end{tikzpicture}
\end{center}

\vspace{0.5em}

\begin{minipage}[t]{0.48\textwidth}
\textbf{Weg 1: Formaler Rechenweg}

\textit{Phase 1 (beschleunigt):}
$$s_1 = \dfrac{1}{2} \cdot a \cdot t_1^2$$
$$s_1 = \dfrac{1}{2} \cdot \SI{5}{\metre\per\second\squared} \cdot (\SI{4}{\second})^2$$
$$s_1 = \dfrac{1}{2} \cdot \SI{5}{\metre\per\second\squared} \cdot \SI{16}{\second\squared} = \SI{40}{\metre}$$

\textit{Phase 2 (gleichförmig):}
$$v_1 = a \cdot t_1 = \SI{5}{\metre\per\second\squared} \cdot \SI{4}{\second} = \SI{20}{\metre\per\second}$$
$$s_2 = v_1 \cdot t_2 = \SI{20}{\metre\per\second} \cdot \SI{4}{\second} = \SI{80}{\metre}$$

\textit{Gesamt:}
$$s_{\text{ges}} = s_1 + s_2 = \SI{40}{\metre} + \SI{80}{\metre} = \SI{120}{\metre}$$
\end{minipage}
\hfill
\begin{minipage}[t]{0.48\textwidth}
\textbf{Weg 2: Grafischer Weg (Flächen)}

\textit{Fläche $A_1$ (Dreieck):}
$$A_1 = \dfrac{1}{2} \cdot \text{Grundseite} \cdot \text{Höhe}$$
$$A_1 = \dfrac{1}{2} \cdot \SI{4}{\second} \cdot \SI{20}{\metre\per\second}$$
$$A_1 = \SI{40}{\metre}$$

\textit{Fläche $A_2$ (Rechteck):}
$$A_2 = \text{Breite} \cdot \text{Höhe}$$
$$A_2 = \SI{4}{\second} \cdot \SI{20}{\metre\per\second}$$
$$A_2 = \SI{80}{\metre}$$

\textit{Gesamt:}
$$s_{\text{ges}} = A_1 + A_2 = \SI{40}{\metre} + \SI{80}{\metre} = \SI{120}{\metre}$$
\end{minipage}

\vspace{0.5em}

\begin{center}
\fbox{\textbf{Beide Wege führen zum gleichen Ergebnis: $s = \SI{120}{\metre}$}}
\end{center}
\end{beispielbox}

\begin{warnbox}
Bei der Flächenberechnung im $v$-$t$-Diagramm:
\begin{itemize}[leftmargin=1em]
    \item \textbf{Rechteck:} gleichförmige Bewegung ($s = v \cdot t$)
    \item \textbf{Dreieck:} Beschleunigung aus dem Stand ($s = \frac{1}{2} \cdot v \cdot t = \frac{1}{2} a t^2$)
    \item \textbf{Trapez:} Beschleunigung mit Anfangsgeschwindigkeit
\end{itemize}
\end{warnbox}

\newpage

% ═══════════════════════════════════════════════════════════════
% KAPITEL 2: GLEICHMÄSSIG BESCHLEUNIGTE BEWEGUNG
% ═══════════════════════════════════════════════════════════════

\section{Gleichmäßig beschleunigte Bewegung}

\subsection{Definition der Beschleunigung}

\begin{merksatz}
Die \wichtig{Beschleunigung $a$} gibt an, wie schnell sich die Geschwindigkeit ändert. Bei einer \wichtig{gleichmäßig beschleunigten Bewegung} ist $a$ konstant.
\end{merksatz}

\begin{formelbox}[title=Beschleunigung]
\begin{center}
{\Large $a = \dfrac{\Delta v}{\Delta t} = \dfrac{v_1 - v_0}{t}$}
\end{center}
\vspace{0.3em}
\begin{tabular}{lll}
$a$ & = & Beschleunigung in $\si{\metre\per\second\squared}$ \\
$\Delta v$ & = & Geschwindigkeitsänderung in $\si{\metre\per\second}$ \\
$\Delta t$ & = & Zeitänderung in $\si{\second}$ \\
\end{tabular}
\end{formelbox}

\begin{warnbox}
\textbf{Verzögerung} = negative Beschleunigung (Abbremsen): $a < 0$
\end{warnbox}

\subsection{Die drei Grundformeln}

Für Bewegung \textbf{aus dem Stand} ($v_0 = 0$, $s_0 = 0$):

\begin{formelbox}[title=Formeln der gleichmäßig beschleunigten Bewegung]
\begin{center}
\renewcommand{\arraystretch}{2.5}
\begin{tabular}{|c|c|c|}
\hline
\textbf{Geschwindigkeit} & \textbf{Weg} & \textbf{Weg (ohne $t$)} \\
\hline
$v = a \cdot t$ & $s = \dfrac{1}{2} a \cdot t^2$ & $s = \dfrac{v^2}{2a}$ \\
\hline
\end{tabular}
\end{center}
\end{formelbox}

\begin{beispielbox}
\textbf{Rechenbeispiel:} Ein Auto beschleunigt mit $a = \SI{4}{\metre\per\second\squared}$ für $t = \SI{5}{\second}$.

\begin{itemize}[leftmargin=1em]
    \item Geschwindigkeit: $v = a \cdot t = \SI{4}{\metre\per\second\squared} \cdot \SI{5}{\second} = \SI{20}{\metre\per\second} = \SI{72}{\kilo\metre\per\hour}$
    \item Weg: $s = \dfrac{1}{2} \cdot a \cdot t^2 = \dfrac{1}{2} \cdot \SI{4}{\metre\per\second\squared} \cdot (\SI{5}{\second})^2 = \SI{50}{\metre}$
\end{itemize}
\end{beispielbox}

\subsection{Diagramme der beschleunigten Bewegung}

\textbf{Beschleunigung ($a > 0$):}

\begin{minipage}{0.48\textwidth}
\centering
\textbf{$v$-$t$-Diagramm (Anfahren)}

\begin{tikzpicture}[scale=0.8]
    \begin{axis}[
        xlabel={$t$ in s},
        ylabel={$v$ in m/s},
        xmin=0, xmax=6,
        ymin=0, ymax=25,
        grid=major,
        width=7cm, height=5cm,
        axis lines=left,
    ]
    \addplot[very thick, physik, domain=0:5] {4*x};
    \fill[physik, opacity=0.2] (0,0) -- (5,0) -- (5,20) -- cycle;
    \end{axis}
\end{tikzpicture}

Steigung = Beschleunigung $a > 0$

Fläche = Weg $s$
\end{minipage}
\hfill
\begin{minipage}{0.48\textwidth}
\centering
\textbf{$s$-$t$-Diagramm (Anfahren)}

\begin{tikzpicture}[scale=0.8]
    \begin{axis}[
        xlabel={$t$ in s},
        ylabel={$s$ in m},
        xmin=0, xmax=6,
        ymin=0, ymax=60,
        grid=major,
        width=7cm, height=5cm,
        axis lines=left,
    ]
    \addplot[very thick, physik, domain=0:5, samples=50] {0.5*4*x^2};
    \end{axis}
\end{tikzpicture}

Parabel (nach oben geöffnet)

Steigung nimmt zu $\rightarrow$ $v$ steigt
\end{minipage}

\vspace{1em}

\textbf{Verzögerung / Bremsen ($a < 0$):}

\begin{minipage}{0.48\textwidth}
\centering
\textbf{$v$-$t$-Diagramm (Bremsen)}

\begin{tikzpicture}[scale=0.8]
    \begin{axis}[
        xlabel={$t$ in s},
        ylabel={$v$ in m/s},
        xmin=0, xmax=6,
        ymin=0, ymax=25,
        grid=major,
        width=7cm, height=5cm,
        axis lines=left,
    ]
    \addplot[very thick, warning, domain=0:5] {20-4*x};
    \fill[warning, opacity=0.2] (0,0) -- (0,20) -- (5,0) -- cycle;
    \end{axis}
\end{tikzpicture}

Steigung = Verzögerung $a < 0$

Fläche = Bremsweg $s$
\end{minipage}
\hfill
\begin{minipage}{0.48\textwidth}
\centering
\textbf{$s$-$t$-Diagramm (Bremsen)}

\begin{tikzpicture}[scale=0.8]
    \begin{axis}[
        xlabel={$t$ in s},
        ylabel={$s$ in m},
        xmin=0, xmax=6,
        ymin=0, ymax=60,
        grid=major,
        width=7cm, height=5cm,
        axis lines=left,
    ]
    \addplot[very thick, warning, domain=0:5, samples=50] {20*x - 0.5*4*x^2};
    \end{axis}
\end{tikzpicture}

Parabel (nach unten geöffnet)

Steigung nimmt ab $\rightarrow$ $v$ sinkt
\end{minipage}

\vspace{1em}

\begin{beispielbox}[title=Alltagsbeispiel: Ampelfahrt]
\textbf{Situation:} Ein Auto startet an einer roten Ampel, fährt dann gleichförmig und bremst vor der nächsten Ampel.

\begin{center}
\begin{tikzpicture}[scale=0.9]
    \begin{axis}[
        xlabel={$t$ in s},
        ylabel={$s$ in m},
        xmin=0, xmax=18,
        ymin=0, ymax=260,
        grid=major,
        width=13cm, height=6cm,
        axis lines=left,
        xtick={0,2,4,6,8,10,12,14,16},
        ytick={0,40,80,120,160,200,240},
        legend pos=north west,
        legend style={font=\small},
    ]
    % Phase 1: Anfahren (0-4s), a=5 m/s², v geht von 0 auf 20 m/s
    \addplot[very thick, physik, domain=0:4, samples=50] {0.5*5*x^2};
    % Phase 2: Gleichförmig (4-12s), v=20 m/s, s_0=40m
    \addplot[very thick, success, domain=4:12] {40 + 20*(x-4)};
    % Phase 3: Bremsen (12-16s), a=-5 m/s², v geht von 20 auf 0
    \addplot[very thick, warning, domain=12:16, samples=50] {200 + 20*(x-12) - 0.5*5*(x-12)^2};

    % Markierungen
    \node[below] at (2, -15) {\small\textcolor{physik}{Anfahren}};
    \node[below] at (8, -15) {\small\textcolor{success}{Gleichförmig}};
    \node[below] at (14, -15) {\small\textcolor{warning}{Bremsen}};
    \end{axis}
\end{tikzpicture}
\end{center}

\textbf{Interpretation des $s$-$t$-Diagramms:}
\begin{itemize}[leftmargin=1em]
    \item \textcolor{physik}{\textbf{Phase 1 (Anfahren):}} Parabel nach oben geöffnet -- Steigung nimmt zu (Geschwindigkeit steigt)
    \item \textcolor{success}{\textbf{Phase 2 (Gleichförmig):}} Gerade -- konstante Steigung (konstante Geschwindigkeit)
    \item \textcolor{warning}{\textbf{Phase 3 (Bremsen):}} Parabel nach unten geöffnet -- Steigung nimmt ab (Geschwindigkeit sinkt)
\end{itemize}

\textbf{Zugehöriges $v$-$t$-Diagramm:}

\begin{center}
\begin{tikzpicture}[scale=0.9]
    \begin{axis}[
        xlabel={$t$ in s},
        ylabel={$v$ in m/s},
        xmin=0, xmax=18,
        ymin=0, ymax=25,
        grid=major,
        width=13cm, height=4cm,
        axis lines=left,
        xtick={0,2,4,6,8,10,12,14,16},
    ]
    % Phase 1: Anfahren
    \addplot[very thick, physik, domain=0:4] {5*x};
    % Phase 2: Gleichförmig
    \addplot[very thick, success, domain=4:12] {20};
    % Phase 3: Bremsen
    \addplot[very thick, warning, domain=12:16] {20 - 5*(x-12)};
    \end{axis}
\end{tikzpicture}
\end{center}
\end{beispielbox}

\newpage

% ═══════════════════════════════════════════════════════════════
% KAPITEL 3: ZUSAMMENGESETZTE BEWEGUNGEN
% ═══════════════════════════════════════════════════════════════

\section{Zusammengesetzte Bewegungen}

\subsection{Prinzip}

\begin{merksatz}
Bei zusammengesetzten Bewegungen werden mehrere Bewegungsphasen kombiniert:
\begin{itemize}[leftmargin=1em]
    \item Jede Phase hat eigene Anfangsbedingungen ($v_0$, $s_0$)
    \item Am Übergang: Endwerte der ersten Phase = Anfangswerte der nächsten
\end{itemize}
\end{merksatz}

\subsection{Bremsvorgang}

\begin{formelbox}[title=Bremsweg und Bremszeit]
Bei Anfangsgeschwindigkeit $v_0$ und Verzögerung $a$ (als Betrag):

\begin{center}
{\large $s_{\text{Brems}} = \dfrac{v_0^2}{2a}$ \hspace{3cm} $t_{\text{Brems}} = \dfrac{v_0}{a}$}
\end{center}
\end{formelbox}

\begin{beispielbox}
\textbf{Bremsbeispiel:} Auto mit $v_0 = \SI{20}{\metre\per\second}$ bremst mit $a = \SI{5}{\metre\per\second\squared}$.

\begin{itemize}[leftmargin=1em]
    \item Bremsweg: $s = \dfrac{v_0^2}{2a} = \dfrac{(\SI{20}{\metre\per\second})^2}{2 \cdot \SI{5}{\metre\per\second\squared}} = \dfrac{\SI{400}{\metre\squared\per\second\squared}}{\SI{10}{\metre\per\second\squared}} = \SI{40}{\metre}$
    \item Bremszeit: $t = \dfrac{v_0}{a} = \dfrac{\SI{20}{\metre\per\second}}{\SI{5}{\metre\per\second\squared}} = \SI{4}{\second}$
\end{itemize}
\end{beispielbox}

\subsection{Bewegung mit Anfangsbedingungen}

Für Bewegung mit \textbf{Anfangsgeschwindigkeit $v_0$} und \textbf{Startposition $s_0$}:

\begin{formelbox}[title=Allgemeine Bewegungsgleichungen]
\begin{center}
{\large $v(t) = v_0 + a \cdot t$}

\vspace{0.5em}

{\large $s(t) = s_0 + v_0 \cdot t + \dfrac{1}{2} a \cdot t^2$}
\end{center}
\end{formelbox}

\subsection{Relative Geschwindigkeiten}

Bei der Betrachtung von zwei bewegten Körpern ist die \wichtig{Relativgeschwindigkeit} entscheidend:

\begin{formelbox}[title=Relativgeschwindigkeit]
\begin{itemize}[leftmargin=1em]
    \item \textbf{Gleiche Richtung:} $v_{\text{rel}} = v_1 - v_2$ (Differenz)
    \item \textbf{Entgegengesetzte Richtung:} $v_{\text{rel}} = v_1 + v_2$ (Summe)
\end{itemize}
\end{formelbox}

\begin{beispielbox}
\textbf{Beispiel 1: Zwei Autos auf der Autobahn}

Auto 1 fährt mit $v_1 = \SI{120}{\kilo\metre\per\hour}$, Auto 2 mit $v_2 = \SI{100}{\kilo\metre\per\hour}$ (gleiche Richtung).

Die Relativgeschwindigkeit beträgt: $v_{\text{rel}} = \SI{120}{\kilo\metre\per\hour} - \SI{100}{\kilo\metre\per\hour} = \SI{20}{\kilo\metre\per\hour}$

$\Rightarrow$ Auto 1 nähert sich Auto 2 mit $\SI{20}{\kilo\metre\per\hour}$.
\end{beispielbox}

\begin{beispielbox}
\textbf{Beispiel 2: Frontalverkehr}

Zwei Autos fahren aufeinander zu, jeweils mit $v = \SI{50}{\kilo\metre\per\hour}$.

Die Relativgeschwindigkeit beträgt: $v_{\text{rel}} = \SI{50}{\kilo\metre\per\hour} + \SI{50}{\kilo\metre\per\hour} = \SI{100}{\kilo\metre\per\hour}$

$\Rightarrow$ Der Abstand verringert sich mit $\SI{100}{\kilo\metre\per\hour}$.
\end{beispielbox}

\subsection{Überholmanöver}

\begin{beispielbox}
\textbf{Typische Aufgabe:} LKW fährt mit $v_{\text{LKW}} = \SI{20}{\metre\per\second}$. Auto startet mit $v_0 = 0$ und beschleunigt mit $a = \SI{4}{\metre\per\second\squared}$. Wann und wo überholt das Auto?

\textbf{Lösungsweg:}

1. \textbf{Gleichungen aufstellen:}
\begin{align*}
    s_{\text{LKW}} &= v_{\text{LKW}} \cdot t = \SI{20}{\metre\per\second} \cdot t \\
    s_{\text{Auto}} &= \dfrac{1}{2} a \cdot t^2 = \dfrac{1}{2} \cdot \SI{4}{\metre\per\second\squared} \cdot t^2 = \SI{2}{\metre\per\second\squared} \cdot t^2
\end{align*}

2. \textbf{Gleichsetzen:} $s_{\text{LKW}} = s_{\text{Auto}}$
$$\SI{20}{\metre\per\second} \cdot t = \SI{2}{\metre\per\second\squared} \cdot t^2$$
$$t \cdot (\SI{20}{\metre\per\second} - \SI{2}{\metre\per\second\squared} \cdot t) = 0 \quad \Rightarrow \quad t = 0 \text{ oder } t = \SI{10}{\second}$$

3. \textbf{Ort berechnen:}
$$s = \SI{20}{\metre\per\second} \cdot \SI{10}{\second} = \SI{200}{\metre}$$

\textbf{Antwort:} Das Auto überholt nach $\SI{10}{\second}$ bei $\SI{200}{\metre}$.
\end{beispielbox}

\newpage

% ═══════════════════════════════════════════════════════════════
% KAPITEL 4: DER FREIE FALL
% ═══════════════════════════════════════════════════════════════

\section{Der freie Fall}

\subsection{Definition}

\begin{merksatz}
Der \wichtig{freie Fall} ist eine gleichmäßig beschleunigte Bewegung unter dem Einfluss der Erdbeschleunigung $g$. Der Luftwiderstand wird vernachlässigt.

\begin{center}
{\Large $g = \SI{9,81}{\metre\per\second\squared}$} \quad (Kopfrechnen: $g \approx \SI{10}{\metre\per\second\squared}$)
\end{center}
\end{merksatz}

\begin{beispielbox}[title=Historischer Hintergrund: Galileo Galilei]
\textbf{Galileo Galilei} (1564--1642) entdeckte durch systematische Experimente, dass alle Körper \textbf{gleich schnell fallen} -- unabhängig von ihrer Masse. Legende: Fallversuche am Schiefen Turm von Pisa.

\textbf{Warum Feder und Stein unterschiedlich fallen:}
Der Unterschied liegt nicht an der Masse, sondern am \textbf{Luftwiderstand}. Im Vakuum fallen Feder und Hammer exakt gleich schnell (Beweis: Apollo-15-Mission auf dem Mond, 1971).
\end{beispielbox}

\subsection{Fallbeschleunigung auf anderen Himmelskörpern}

\begin{center}
\begin{tabular}{|l|c|c|}
\hline
\textbf{Himmelskörper} & \textbf{$g$ in m/s²} & \textbf{Vergleich zur Erde} \\
\hline
Mond & 1,62 & $\approx 1/6$ \\
Mars & 3,71 & $\approx 1/3$ \\
Jupiter & 24,79 & $\approx 2,5\times$ \\
Sonne & 274 & $\approx 28\times$ \\
\hline
\end{tabular}
\end{center}

\begin{beispielbox}[title=Mondmission]
Ein Astronaut lässt auf dem Mond ($g_{\text{Mond}} = \SI{1,6}{\metre\per\second\squared}$) einen Hammer aus $\SI{1,8}{\metre}$ Höhe fallen.

\textbf{Fallzeit:} $t = \sqrt{\dfrac{2h}{g}} = \sqrt{\dfrac{2 \cdot \SI{1,8}{\metre}}{\SI{1,6}{\metre\per\second\squared}}} = \sqrt{\SI{2,25}{\second\squared}} = \SI{1,5}{\second}$

\textbf{Vergleich Erde:} $t = \sqrt{\dfrac{2 \cdot \SI{1,8}{\metre}}{\SI{10}{\metre\per\second\squared}}} = \sqrt{\SI{0,36}{\second\squared}} = \SI{0,6}{\second}$

$\Rightarrow$ Auf dem Mond fällt der Hammer $\approx 2,5$-mal langsamer!
\end{beispielbox}

\subsection{Formeln des freien Falls}

\begin{formelbox}[title=Formeln des freien Falls (Start aus Ruhe)]
\begin{center}
\renewcommand{\arraystretch}{2.2}
\begin{tabular}{|c|c|c|}
\hline
\textbf{Geschwindigkeit} & \textbf{Fallhöhe} & \textbf{Fallhöhe (ohne $t$)} \\
\hline
$v = g \cdot t$ & $h = \dfrac{1}{2} g \cdot t^2$ & $h = \dfrac{v^2}{2g}$ \\
\hline
\end{tabular}

\vspace{1em}

\textbf{Umgestellt:}

\begin{tabular}{|c|c|}
\hline
\textbf{Fallzeit} & \textbf{Aufprallgeschwindigkeit} \\
\hline
$t = \sqrt{\dfrac{2h}{g}}$ & $v = \sqrt{2gh}$ \\
\hline
\end{tabular}
\end{center}
\end{formelbox}

\subsection{Wertetabelle (Kopfrechnen mit $g \approx \SI{10}{\metre\per\second\squared}$)}

\begin{center}
\begin{tabular}{|c|c|c|c|c|c|c|}
\hline
$t$ in s & 0 & 1 & 2 & 3 & 4 & 5 \\
\hline
$v$ in m/s & 0 & 10 & 20 & 30 & 40 & 50 \\
\hline
$h$ in m & 0 & 5 & 20 & 45 & 80 & 125 \\
\hline
\end{tabular}
\end{center}

\subsection{Diagramme des freien Falls}

\begin{minipage}{0.48\textwidth}
\centering
\textbf{$v$-$t$-Diagramm}

\begin{tikzpicture}[scale=0.75]
    \begin{axis}[
        xlabel={$t$ in s},
        ylabel={$v$ in m/s},
        xmin=0, xmax=5,
        ymin=0, ymax=55,
        grid=major,
        width=6.5cm, height=5cm,
        axis lines=left,
    ]
    \addplot[very thick, physik, domain=0:5] {10*x};
    \end{axis}
\end{tikzpicture}

Gerade mit Steigung $g$
\end{minipage}
\hfill
\begin{minipage}{0.48\textwidth}
\centering
\textbf{$h$-$t$-Diagramm}

\begin{tikzpicture}[scale=0.75]
    \begin{axis}[
        xlabel={$t$ in s},
        ylabel={$h$ in m},
        xmin=0, xmax=5,
        ymin=0, ymax=140,
        grid=major,
        width=6.5cm, height=5cm,
        axis lines=left,
    ]
    \addplot[very thick, physik, domain=0:5, samples=50] {5*x^2};
    \end{axis}
\end{tikzpicture}

Parabel: $h = 5t^2$
\end{minipage}

\subsection{Anwendung: Reaktionszeit-Experiment}

\begin{beispielbox}[title=Lineal-Experiment zur Reaktionszeit]
\textbf{Durchführung:} Eine Person hält ein Lineal ($\SI{30}{\centi\metre}$) am oberen Ende. Eine zweite Person hält die Finger bereit (ohne das Lineal zu berühren) an der 0-cm-Marke. Das Lineal wird losgelassen -- die zweite Person fängt es so schnell wie möglich.

\textbf{Auswertung:} Die Fallstrecke $h$ (abgelesen am Lineal) ermöglicht die Berechnung der Reaktionszeit:

$$t_{\text{Reaktion}} = \sqrt{\dfrac{2h}{g}}$$

\textbf{Beispiel:} Das Lineal fällt $h = \SI{20}{\centi\metre} = \SI{0,2}{\metre}$

$$t = \sqrt{\dfrac{2 \cdot \SI{0,2}{\metre}}{\SI{10}{\metre\per\second\squared}}} = \sqrt{\SI{0,04}{\second\squared}} = \SI{0,2}{\second} = \SI{200}{\milli\second}$$

\textbf{Typische Werte:}
\begin{itemize}[leftmargin=1em]
    \item Sehr schnell: $< \SI{150}{\milli\second}$ (Fallstrecke $< \SI{11}{\centi\metre}$)
    \item Durchschnitt: $\SI{200}{\milli\second}$ -- $\SI{250}{\milli\second}$ (Fallstrecke $\SI{20}{\centi\metre}$ -- $\SI{31}{\centi\metre}$)
    \item Langsam: $> \SI{300}{\milli\second}$ (Fallstrecke $> \SI{45}{\centi\metre}$)
\end{itemize}
\end{beispielbox}

\subsection{Senkrechter Wurf nach oben}

\begin{formelbox}[title=Senkrechter Wurf nach oben]
Mit Anfangsgeschwindigkeit $v_0$ nach oben:

\begin{center}
$v(t) = v_0 - g \cdot t$ \hspace{2cm} $h(t) = v_0 \cdot t - \dfrac{1}{2} g \cdot t^2$
\end{center}

\textbf{Steigzeit:} $t_{\text{steig}} = \dfrac{v_0}{g}$ \hspace{2cm} \textbf{Maximale Höhe:} $h_{\text{max}} = \dfrac{v_0^2}{2g}$

Am höchsten Punkt gilt: $v = 0$
\end{formelbox}

\begin{merksatz}[title=Symmetrie des senkrechten Wurfs]
Der senkrechte Wurf nach oben ist \textbf{symmetrisch}:
\begin{itemize}[leftmargin=1em]
    \item \textbf{Steigzeit = Fallzeit:} $t_{\text{steig}} = t_{\text{fall}} = \dfrac{v_0}{g}$
    \item \textbf{Gesamtflugzeit:} $t_{\text{ges}} = 2 \cdot t_{\text{steig}} = \dfrac{2 \cdot v_0}{g}$
    \item \textbf{Aufprallgeschwindigkeit = Abwurfgeschwindigkeit:} $v_{\text{Aufprall}} = v_0$
\end{itemize}
\end{merksatz}

\begin{center}
\begin{tikzpicture}[scale=0.8]
    \begin{axis}[
        xlabel={$t$ in s},
        ylabel={$h$ in m},
        xmin=0, xmax=5,
        ymin=0, ymax=25,
        grid=major,
        width=10cm, height=6cm,
        axis lines=left,
        xtick={0,1,2,3,4},
        ytick={0,5,10,15,20},
    ]
    % h(t) = 20t - 5t² für v_0 = 20 m/s
    \addplot[very thick, physik, domain=0:4, samples=50] {20*x - 5*x^2};
    % Markierungen
    \node[above] at (2,20) {$h_{\text{max}}$};
    \draw[dashed, grau] (2,0) -- (2,20);
    \node[below] at (2,0) {$t_{\text{steig}}$};
    \node[below] at (4,0) {$t_{\text{ges}}$};
    \end{axis}
\end{tikzpicture}
\end{center}

\textbf{Interpretation:} Steigphase (links vom Maximum) und Fallphase (rechts) sind spiegelsymmetrisch. Die Parabel ist achsensymmetrisch zur Senkrechten durch den Scheitelpunkt.

\newpage

% ═══════════════════════════════════════════════════════════════
% KAPITEL 5: LUFTWIDERSTAND UND GRENZGESCHWINDIGKEIT
% ═══════════════════════════════════════════════════════════════

\section{Luftwiderstand und Grenzgeschwindigkeit}

\begin{warnbox}[title=Hinweis zur Klausur]
Dieses Thema wird \textbf{nur qualitativ} (ohne Rechnung) geprüft. Du solltest:
\begin{itemize}[leftmargin=1em]
    \item Die \textbf{Begriffe} kennen (Luftwiderstand, Grenzgeschwindigkeit)
    \item Die \textbf{Ursache} der Grenzgeschwindigkeit erklären können ($F_W = F_G$)
    \item Das \textbf{$v$-$t$-Diagramm} eines Fallschirmspringers interpretieren können
\end{itemize}
\end{warnbox}

\subsection{Die Luftwiderstandskraft}

\begin{formelbox}[title=Luftwiderstandskraft]
\begin{center}
{\Large $F_W = \dfrac{1}{2} \cdot \rho \cdot c_W \cdot A \cdot v^2$}
\end{center}
\vspace{0.3em}
\begin{tabular}{lll}
$F_W$ & = & Luftwiderstandskraft in $\si{\newton}$ \\
$\rho$ & = & Luftdichte ($\approx \SI{1,2}{\kilo\gram\per\metre\cubed}$) \\
$c_W$ & = & Widerstandsbeiwert (dimensionslos) \\
$A$ & = & Stirnfläche (Querschnittsfläche) in $\si{\metre\squared}$ \\
$v$ & = & Geschwindigkeit in $\si{\metre\per\second}$ \\
\end{tabular}
\end{formelbox}

\begin{merksatz}
Die Luftwiderstandskraft wächst \wichtig{quadratisch} mit der Geschwindigkeit: Doppelte Geschwindigkeit $\Rightarrow$ vierfacher Luftwiderstand!
\end{merksatz}

\subsection{Der Widerstandsbeiwert $c_W$}

\begin{center}
\begin{tabular}{|l|c|}
\hline
\textbf{Körperform} & \textbf{$c_W$-Wert} \\
\hline
Tropfenform (optimal) & 0,04 \\
Sportwagen & 0,25--0,35 \\
PKW (normal) & 0,30--0,40 \\
LKW & 0,60--0,90 \\
Motorradfahrer & 0,80--1,00 \\
Fallschirmspringer (gestreckt) & 0,90 \\
Fallschirm (offen) & 1,30--1,50 \\
\hline
\end{tabular}
\end{center}

\subsection{Die Grenzgeschwindigkeit}

\begin{merksatz}
Die \wichtig{Grenzgeschwindigkeit} $v_{\text{max}}$ wird erreicht, wenn der Luftwiderstand gleich der Gewichtskraft ist: $F_W = F_G$

Der Körper fällt dann mit \textbf{konstanter Geschwindigkeit} (keine weitere Beschleunigung).
\end{merksatz}

\begin{formelbox}[title=Grenzgeschwindigkeit]
Aus $F_W = F_G$ folgt:
\begin{center}
{\Large $v_{\text{max}} = \sqrt{\dfrac{2 \cdot m \cdot g}{\rho \cdot c_W \cdot A}}$}
\end{center}
\end{formelbox}

\subsubsection*{Herleitung der Formel (zur Information)}

Kräftegleichgewicht: $F_W = F_G$

$$\dfrac{1}{2} \cdot \rho \cdot c_W \cdot A \cdot v^2 = m \cdot g$$

Nach $v$ umstellen:

$$v^2 = \dfrac{2 \cdot m \cdot g}{\rho \cdot c_W \cdot A} \quad \Rightarrow \quad v_{\text{max}} = \sqrt{\dfrac{2 \cdot m \cdot g}{\rho \cdot c_W \cdot A}}$$

\subsubsection*{Parameteranalyse: Was beeinflusst $v_{\text{max}}$?}

\begin{center}
\begin{tabular}{|l|c|l|}
\hline
\textbf{Parameter} & \textbf{Änderung} & \textbf{Effekt auf $v_{\text{max}}$} \\
\hline
Masse $m$ & $\uparrow$ verdoppelt & $v_{\text{max}} \cdot \sqrt{2} \approx 1,4 \times$ \\
\hline
Stirnfläche $A$ & $\uparrow$ verdoppelt & $v_{\text{max}} / \sqrt{2} \approx 0,7 \times$ \\
\hline
$c_W$-Wert & $\uparrow$ verdoppelt & $v_{\text{max}} / \sqrt{2} \approx 0,7 \times$ \\
\hline
Luftdichte $\rho$ & $\downarrow$ (große Höhe) & $v_{\text{max}}$ steigt \\
\hline
\end{tabular}
\end{center}

\begin{warnbox}
\textbf{Wichtig:} Schwere Objekte erreichen höhere Grenzgeschwindigkeiten! Ein fallender Elefant fällt schneller als eine fallende Maus (bei gleicher Körperform).
\end{warnbox}

\begin{beispielbox}[title=Red Bull Stratos (2012)]
Felix Baumgartner sprang aus $\SI{39}{\kilo\metre}$ Höhe -- fast aus dem Weltraum!

\textbf{Besonderheit:} In großer Höhe ist die Luftdichte $\rho$ sehr gering. Daher erreichte er im freien Fall eine Geschwindigkeit von $v_{\text{max}} \approx \SI{1358}{\kilo\metre\per\hour}$ -- schneller als der Schall!

Je tiefer er fiel, desto dichter wurde die Luft, und er bremste automatisch auf $\approx \SI{200}{\kilo\metre\per\hour}$ ab, bevor er den Fallschirm öffnete.
\end{beispielbox}

\subsection{Medienvergleich: Luft vs. Wasser}

\begin{center}
\begin{tabular}{|l|c|c|}
\hline
\textbf{Medium} & \textbf{Dichte $\rho$} & \textbf{Widerstand} \\
\hline
Luft (Meereshöhe) & $\SI{1,2}{\kilo\gram\per\metre\cubed}$ & gering \\
Wasser & $\SI{1000}{\kilo\gram\per\metre\cubed}$ & sehr hoch \\
\hline
\end{tabular}
\end{center}

\begin{merksatz}
Wasser ist $\approx 800$-mal dichter als Luft. Daher ist der Strömungswiderstand in Wasser deutlich größer. Schwimmer erreichen viel niedrigere Grenzgeschwindigkeiten als Fallschirmspringer.
\end{merksatz}

\begin{beispielbox}
\textbf{Fallschirmspringer:}
\begin{itemize}[leftmargin=1em]
    \item Ohne Fallschirm: $v_{\text{max}} \approx \SI{50}{\metre\per\second} = \SI{180}{\kilo\metre\per\hour}$
    \item Mit Fallschirm: $v_{\text{max}} \approx \SI{5}{\metre\per\second} = \SI{18}{\kilo\metre\per\hour}$
\end{itemize}
\end{beispielbox}

\textbf{$v$-$t$-Diagramm} (Geschwindigkeit über Zeit):

\begin{center}
\begin{tikzpicture}[scale=0.85]
    \begin{axis}[
        xlabel={$t$ in s},
        ylabel={$v$ in m/s},
        xmin=0, xmax=20,
        ymin=0, ymax=55,
        grid=major,
        width=11cm, height=5cm,
        axis lines=left,
        legend pos=south east,
    ]
    % Freier Fall ohne Luftwiderstand
    \addplot[very thick, physik, dashed, domain=0:5] {10*x};
    \addlegendentry{ohne Luftwiderstand}
    % Mit Luftwiderstand (annähernd)
    \addplot[very thick, success, domain=0:20, samples=50] {50*(1-exp(-0.2*x))};
    \addlegendentry{mit Luftwiderstand}
    % Grenzgeschwindigkeit
    \draw[thick, warning, dashed] (0,50) -- (20,50) node[right] {$v_{\text{max}}$};
    \end{axis}
\end{tikzpicture}
\end{center}

\vspace{0.5em}

\textbf{$s$-$t$-Diagramm} (Fallstrecke über Zeit):

\begin{center}
\begin{tikzpicture}[scale=0.85]
    \begin{axis}[
        xlabel={$t$ in s},
        ylabel={$s$ in m},
        xmin=0, xmax=20,
        ymin=0, ymax=1100,
        grid=major,
        width=11cm, height=5cm,
        axis lines=left,
        legend pos=north west,
    ]
    % Freier Fall ohne Luftwiderstand (Parabel)
    \addplot[very thick, physik, dashed, domain=0:14, samples=50] {5*x^2};
    \addlegendentry{ohne Luftwiderstand}
    % Mit Luftwiderstand: erst Parabel, dann linear
    \addplot[very thick, success, domain=0:20, samples=100] {50*x + 250*(exp(-0.2*x)-1)};
    \addlegendentry{mit Luftwiderstand}
    \end{axis}
\end{tikzpicture}
\end{center}

\begin{merksatz}
\textbf{Interpretation:} Bei Grenzgeschwindigkeit wird das $s$-$t$-Diagramm zur \textbf{Geraden} (konstante Steigung = konstante Geschwindigkeit). Das $v$-$t$-Diagramm nähert sich asymptotisch der horizontalen Linie $v_{\text{max}}$.
\end{merksatz}

\newpage

% ═══════════════════════════════════════════════════════════════
% KAPITEL 6: FORMELSAMMLUNG
% ═══════════════════════════════════════════════════════════════

\section{Formelsammlung}

\subsection{Alle wichtigen Formeln auf einen Blick}

\begin{formelbox}[title=Gleichförmige Bewegung]
\begin{center}
\renewcommand{\arraystretch}{1.8}
\begin{tabular}{|c|c|c|}
\hline
$v = \dfrac{s}{t}$ & $s = v \cdot t$ & $t = \dfrac{s}{v}$ \\
\hline
\end{tabular}
\end{center}
\end{formelbox}

\begin{formelbox}[title=Gleichmäßig beschleunigte Bewegung (aus dem Stand)]
\begin{center}
\renewcommand{\arraystretch}{1.8}
\begin{tabular}{|c|c|c|c|}
\hline
$a = \dfrac{\Delta v}{\Delta t}$ & $v = a \cdot t$ & $s = \dfrac{1}{2} a \cdot t^2$ & $s = \dfrac{v^2}{2a}$ \\
\hline
\end{tabular}
\end{center}
\end{formelbox}

\begin{formelbox}[title={Mit Anfangsbedingungen ($v_0 \neq 0$, $s_0 \neq 0$)}]
\begin{center}
\renewcommand{\arraystretch}{1.8}
\begin{tabular}{|c|c|}
\hline
$v(t) = v_0 + a \cdot t$ & $s(t) = s_0 + v_0 \cdot t + \dfrac{1}{2} a \cdot t^2$ \\
\hline
\end{tabular}
\end{center}
\end{formelbox}

\begin{formelbox}[title=Freier Fall und senkrechter Wurf]
\begin{center}
\renewcommand{\arraystretch}{1.8}
\begin{tabular}{|c|c|c|}
\hline
$v = g \cdot t$ & $h = \dfrac{1}{2} g \cdot t^2$ & $v = \sqrt{2gh}$ \\
\hline
$t = \sqrt{\dfrac{2h}{g}}$ & $t_{\text{steig}} = \dfrac{v_0}{g}$ & $h_{\text{max}} = \dfrac{v_0^2}{2g}$ \\
\hline
\end{tabular}
\end{center}
\end{formelbox}

\begin{formelbox}[title=Luftwiderstand (zur Information)]
\begin{center}
\renewcommand{\arraystretch}{1.8}
\begin{tabular}{|c|c|}
\hline
$F_W = \dfrac{1}{2} \rho \cdot c_W \cdot A \cdot v^2$ & $v_{\text{max}} = \sqrt{\dfrac{2mg}{\rho \cdot c_W \cdot A}}$ \\
\hline
\end{tabular}
\end{center}
\end{formelbox}

\subsection{Konstanten und Umrechnungen}

\begin{center}
\begin{tabular}{|l|l|l|}
\hline
\textbf{Konstante} & \textbf{Wert} & \textbf{Kopfrechnen} \\
\hline
Erdbeschleunigung $g$ & $\SI{9,81}{\metre\per\second\squared}$ & $\approx \SI{10}{\metre\per\second\squared}$ \\
Luftdichte $\rho$ & $\SI{1,2}{\kilo\gram\per\metre\cubed}$ & -- \\
\hline
\end{tabular}
\hspace{1cm}
\begin{tabular}{|l|}
\hline
\textbf{Umrechnung} \\
\hline
$\SI{1}{\kilo\metre\per\hour} = \dfrac{1}{3{,}6}\,\si{\metre\per\second}$ \\[0.5em]
$\SI{1}{\metre\per\second} = \SI{3,6}{\kilo\metre\per\hour}$ \\
\hline
\end{tabular}
\end{center}

\subsection{Kopfrechenbare Werte}

\begin{center}
\begin{tabular}{|c|c|c|c|c|c|c|c|c|}
\hline
\textbf{km/h} & 36 & 72 & 90 & 108 & 144 & 180 & 216 & 360 \\
\hline
\textbf{m/s} & 10 & 20 & 25 & 30 & 40 & 50 & 60 & 100 \\
\hline
\end{tabular}
\end{center}

\begin{center}
\begin{tabular}{|c|c|c|c|c|c|c|c|c|c|}
\hline
$\sqrt{x}$ & $\sqrt{4}$ & $\sqrt{9}$ & $\sqrt{16}$ & $\sqrt{25}$ & $\sqrt{36}$ & $\sqrt{49}$ & $\sqrt{64}$ & $\sqrt{81}$ & $\sqrt{100}$ \\
\hline
Ergebnis & 2 & 3 & 4 & 5 & 6 & 7 & 8 & 9 & 10 \\
\hline
\end{tabular}
\end{center}

\newpage

% ═══════════════════════════════════════════════════════════════
% KAPITEL 7: ÜBUNGSAUFGABEN MIT LÖSUNGEN
% ═══════════════════════════════════════════════════════════════

\section{Übungsaufgaben mit Lösungen}

\textbf{Aufgabe 1: Gleichförmige Bewegung}

Ein ICE fährt mit konstanter Geschwindigkeit von $\SI{72}{\kilo\metre\per\hour}$.

a) Rechne die Geschwindigkeit in $\si{\metre\per\second}$ um.

b) Welche Strecke legt der ICE in $\SI{30}{\second}$ zurück?

\begin{tcolorbox}[colback=hellgrau, colframe=grau, boxrule=0.5pt, title=Lösung]
a) $v = \dfrac{\SI{72}{\kilo\metre\per\hour}}{3{,}6} = \SI{20}{\metre\per\second}$

b) $s = v \cdot t = \SI{20}{\metre\per\second} \cdot \SI{30}{\second} = \SI{600}{\metre}$
\end{tcolorbox}

\vspace{0.5em}

\textbf{Aufgabe 2: Beschleunigung}

Ein Sportwagen beschleunigt aus dem Stand in $\SI{5}{\second}$ auf $\SI{108}{\kilo\metre\per\hour}$.

a) Berechne die Beschleunigung.

b) Berechne den zurückgelegten Weg.

\begin{tcolorbox}[colback=hellgrau, colframe=grau, boxrule=0.5pt, title=Lösung]
a) $v_1 = \SI{108}{\kilo\metre\per\hour} = \SI{30}{\metre\per\second}$, \quad $v_0 = 0$

$a = \dfrac{v_1 - v_0}{t} = \dfrac{\SI{30}{\metre\per\second} - 0}{\SI{5}{\second}} = \SI{6}{\metre\per\second\squared}$

b) $s = \dfrac{1}{2} \cdot a \cdot t^2 = \dfrac{1}{2} \cdot \SI{6}{\metre\per\second\squared} \cdot (\SI{5}{\second})^2 = \SI{75}{\metre}$
\end{tcolorbox}

\vspace{0.5em}

\textbf{Aufgabe 3: Bremsvorgang}

Ein Auto fährt mit $\SI{90}{\kilo\metre\per\hour}$ und bremst mit $a = \SI{5}{\metre\per\second\squared}$.

a) Berechne den Bremsweg.

b) Berechne die Bremszeit.

\begin{tcolorbox}[colback=hellgrau, colframe=grau, boxrule=0.5pt, title=Lösung]
$v_0 = \SI{90}{\kilo\metre\per\hour} = \SI{25}{\metre\per\second}$

a) $s = \dfrac{v_0^2}{2a} = \dfrac{(\SI{25}{\metre\per\second})^2}{2 \cdot \SI{5}{\metre\per\second\squared}} = \dfrac{\SI{625}{\metre\squared\per\second\squared}}{\SI{10}{\metre\per\second\squared}} = \SI{62,5}{\metre}$

b) $t = \dfrac{v_0}{a} = \dfrac{\SI{25}{\metre\per\second}}{\SI{5}{\metre\per\second\squared}} = \SI{5}{\second}$
\end{tcolorbox}

\vspace{0.5em}

\textbf{Aufgabe 4: Freier Fall}

Ein Stein fällt von einem $\SI{80}{\metre}$ hohen Turm.

a) Berechne die Fallzeit.

b) Berechne die Aufprallgeschwindigkeit.

\begin{tcolorbox}[colback=hellgrau, colframe=grau, boxrule=0.5pt, title=Lösung (Kopfrechnen mit $g \approx \SI{10}{\metre\per\second\squared}$)]
a) $t = \sqrt{\dfrac{2h}{g}} = \sqrt{\dfrac{2 \cdot \SI{80}{\metre}}{\SI{10}{\metre\per\second\squared}}} = \sqrt{\SI{16}{\second\squared}} = \SI{4}{\second}$

b) $v = g \cdot t = \SI{10}{\metre\per\second\squared} \cdot \SI{4}{\second} = \SI{40}{\metre\per\second} = \SI{144}{\kilo\metre\per\hour}$

\textit{Alternativ:} $v = \sqrt{2gh} = \sqrt{2 \cdot \SI{10}{\metre\per\second\squared} \cdot \SI{80}{\metre}} = \sqrt{\SI{1600}{\metre\squared\per\second\squared}} = \SI{40}{\metre\per\second}$
\end{tcolorbox}

\vspace{0.5em}

\textbf{Aufgabe 5: Senkrechter Wurf}

Ein Ball wird mit $v_0 = \SI{20}{\metre\per\second}$ senkrecht nach oben geworfen.

a) Berechne die maximale Höhe.

b) Berechne die Steigzeit.

\begin{tcolorbox}[colback=hellgrau, colframe=grau, boxrule=0.5pt, title=Lösung (Kopfrechnen mit $g \approx \SI{10}{\metre\per\second\squared}$)]
a) $h_{\text{max}} = \dfrac{v_0^2}{2g} = \dfrac{(\SI{20}{\metre\per\second})^2}{2 \cdot \SI{10}{\metre\per\second\squared}} = \dfrac{\SI{400}{\metre\squared\per\second\squared}}{\SI{20}{\metre\per\second\squared}} = \SI{20}{\metre}$

b) $t_{\text{steig}} = \dfrac{v_0}{g} = \dfrac{\SI{20}{\metre\per\second}}{\SI{10}{\metre\per\second\squared}} = \SI{2}{\second}$
\end{tcolorbox}

\newpage

% ═══════════════════════════════════════════════════════════════
% KAPITEL 8: SELBSTTEST
% ═══════════════════════════════════════════════════════════════

\section{Selbsttest: Habe ich alles verstanden?}

\subsection{Lückentext -- Fülle die Lücken aus!}

\begin{enumerate}
    \item Bei einer gleichförmigen Bewegung ist die \luecke{3cm} konstant.

    \item Die Beschleunigung gibt an, wie schnell sich die \luecke{3cm} ändert.

    \item Im $s$-$t$-Diagramm entspricht die Steigung der \luecke{3cm}.

    \item Im $v$-$t$-Diagramm entspricht die Fläche unter dem Graphen dem \luecke{2cm}.

    \item Die Erdbeschleunigung beträgt $g = $ \luecke{1.5cm} $\si{\metre\per\second\squared}$ (exakt: \luecke{1.5cm} $\si{\metre\per\second\squared}$).

    \item Bei doppelter Geschwindigkeit ist die Luftwiderstandskraft \luecke{2cm} so groß.

    \item Die Grenzgeschwindigkeit wird erreicht, wenn $F_W = $ \luecke{2cm}.

    \item Um km/h in m/s umzurechnen, muss man durch \luecke{1.5cm} dividieren.

    \item Bei doppelter Ausgangsgeschwindigkeit ist der Bremsweg \luecke{2cm} so groß.
\end{enumerate}

\textit{Lösungen: 1) Geschwindigkeit, 2) Geschwindigkeit, 3) Geschwindigkeit, 4) Weg, 5) $\approx$10 / 9,81, 6) viermal, 7) $F_G$, 8) 3,6, 9) viermal}

\subsection{Quick-Check: Richtig oder Falsch?}

\begin{center}
\begin{tabular}{|p{10cm}|c|c|}
\hline
\textbf{Aussage} & \textbf{R} & \textbf{F} \\
\hline
Bei gleichförmiger Bewegung ist die Beschleunigung null. & $\square$ & $\square$ \\
\hline
Im freien Fall ist die Geschwindigkeit konstant. & $\square$ & $\square$ \\
\hline
Der Bremsweg verdoppelt sich bei doppelter Geschwindigkeit. & $\square$ & $\square$ \\
\hline
Die Steigzeit beim senkrechten Wurf ist gleich der Fallzeit. & $\square$ & $\square$ \\
\hline
Ein Fallschirm vergrößert den $c_W$-Wert und die Stirnfläche. & $\square$ & $\square$ \\
\hline
Bei Grenzgeschwindigkeit ist die Beschleunigung null. & $\square$ & $\square$ \\
\hline
$\SI{72}{\kilo\metre\per\hour} = \SI{20}{\metre\per\second}$ & $\square$ & $\square$ \\
\hline
Die Luftwiderstandskraft wächst linear mit der Geschwindigkeit. & $\square$ & $\square$ \\
\hline
\end{tabular}
\end{center}

\textit{Lösungen: R, F, F (vervierfacht!), R, R, R, R, F (quadratisch!)}

\newpage

% ═══════════════════════════════════════════════════════════════
% KAPITEL 9: DIFFERENZIERTE ÜBUNGSAUFGABEN
% ═══════════════════════════════════════════════════════════════

\section{Differenzierte Übungsaufgaben}

\textbf{Niveau I -- Grundlagen} \hfill \textit{[Reproduktion]}

\begin{enumerate}
    \item Rechne um: a) $\SI{54}{\kilo\metre\per\hour}$ in m/s \hspace{1cm} b) $\SI{15}{\metre\per\second}$ in km/h

    \item Ein Radfahrer fährt $\SI{12}{\kilo\metre}$ in $\SI{30}{\minute}$. Berechne seine Geschwindigkeit in km/h und m/s.

    \item Ein Körper fällt $\SI{3}{\second}$ lang. Berechne Fallhöhe und Endgeschwindigkeit.
\end{enumerate}

\vspace{1em}

\textbf{Niveau II -- Anwendung} \hfill \textit{[Transfer]}

\begin{enumerate}[resume]
    \item Ein Auto beschleunigt in $\SI{8}{\second}$ von $\SI{36}{\kilo\metre\per\hour}$ auf $\SI{108}{\kilo\metre\per\hour}$.

    a) Berechne die Beschleunigung.

    b) Berechne den dabei zurückgelegten Weg.

    \item Ein Stein wird mit $\SI{30}{\metre\per\second}$ senkrecht nach oben geworfen.

    a) Berechne die maximale Höhe.

    b) Berechne die gesamte Flugzeit (bis zur Rückkehr).

    \item Bei welcher Höhe hat ein fallender Körper die Geschwindigkeit $\SI{30}{\metre\per\second}$ erreicht?
\end{enumerate}

\vspace{1em}

\textbf{Niveau III -- Vertiefung} \hfill \textit{[Problemlösen]}

\begin{enumerate}[resume]
    \item \textbf{Anhalteweg:} Ein Autofahrer fährt mit $\SI{72}{\kilo\metre\per\hour}$. Die Reaktionszeit beträgt $\SI{1}{\second}$, die Verzögerung $\SI{8}{\metre\per\second\squared}$.

    a) Berechne den Reaktionsweg.

    b) Berechne den Bremsweg.

    c) Berechne den gesamten Anhalteweg.

    \item \textbf{Überholvorgang:} Ein LKW fährt mit $\SI{72}{\kilo\metre\per\hour}$. Ein PKW startet $\SI{100}{\metre}$ dahinter mit $\SI{90}{\kilo\metre\per\hour}$.

    a) Wann hat der PKW den LKW eingeholt?

    b) Welche Strecke hat der PKW dann zurückgelegt?

    \item \textbf{Fallschirmsprung:} Ein Fallschirmspringer ($m = \SI{80}{\kilo\gram}$) erreicht ohne Fallschirm eine Grenzgeschwindigkeit von $\SI{50}{\metre\per\second}$.

    a) Berechne die Luftwiderstandskraft bei Grenzgeschwindigkeit.

    b) Um welchen Faktor muss die Stirnfläche durch den Fallschirm vergrößert werden, damit $v_{\text{max}} = \SI{5}{\metre\per\second}$ beträgt?
\end{enumerate}

\vspace{1em}

\textbf{Diagramm-Übungen (Interpretation und Zeichnen)} \hfill \textit{[AFB II/III]}

\begin{enumerate}[resume]
    \item \textbf{Diagramm interpretieren:} Das folgende $v$-$t$-Diagramm zeigt die Bewegung eines Radfahrers.

\begin{center}
\begin{tikzpicture}[scale=0.7]
    \begin{axis}[
        xlabel={$t$ in s},
        ylabel={$v$ in m/s},
        xmin=0, xmax=25,
        ymin=0, ymax=12,
        grid=major,
        width=10cm, height=5cm,
        axis lines=left,
        xtick={0,5,10,15,20,25},
    ]
    % Phase 1: Beschleunigung (0-5s)
    \addplot[very thick, physik, domain=0:5] {2*x};
    % Phase 2: Gleichförmig (5-15s)
    \addplot[very thick, physik, domain=5:15] {10};
    % Phase 3: Bremsen (15-20s)
    \addplot[very thick, physik, domain=15:20] {10 - 2*(x-15)};
    % Phase 4: Stillstand (20-25s)
    \addplot[very thick, physik, domain=20:25] {0};
    \end{axis}
\end{tikzpicture}
\end{center}

    a) Beschreibe die Bewegung in Worten (vier Phasen).

    b) Berechne die Beschleunigung in Phase 1.

    c) Berechne den gesamten zurückgelegten Weg (Fläche!).

    d) Zeichne das zugehörige $s$-$t$-Diagramm.

    \item \textbf{$s$-$t$-Diagramm zeichnen:} Ein Auto fährt folgende Strecke:
    \begin{itemize}[leftmargin=1.5em]
        \item $t = 0$--$\SI{4}{\second}$: Beschleunigung mit $a = \SI{5}{\metre\per\second\squared}$ aus dem Stand
        \item $t = \SI{4}{\second}$--$\SI{10}{\second}$: Gleichförmige Fahrt
        \item $t = \SI{10}{\second}$--$\SI{14}{\second}$: Bremsen bis zum Stillstand
    \end{itemize}

    a) Berechne die Endgeschwindigkeit nach Phase 1.

    b) Berechne den Weg nach jeder Phase.

    c) Zeichne das $s$-$t$-Diagramm mit korrekter Beschriftung.

    d) Zeichne das $v$-$t$-Diagramm.

    \item \textbf{Mehrphasen-Bewegung analysieren:} Ein Ball wird senkrecht nach oben geworfen und fällt zurück.

\begin{center}
\begin{tikzpicture}[scale=0.7]
    \begin{axis}[
        xlabel={$t$ in s},
        ylabel={$v$ in m/s},
        xmin=0, xmax=5,
        ymin=-25, ymax=25,
        grid=major,
        width=10cm, height=6cm,
        axis lines=center,
        xtick={0,1,2,3,4},
        ytick={-20,-10,0,10,20},
    ]
    % v(t) = 20 - 10t
    \addplot[very thick, physik, domain=0:4] {20 - 10*x};
    \draw[dashed, grau] (0,0) -- (4,0);
    \end{axis}
\end{tikzpicture}
\end{center}

    a) Lies die Anfangsgeschwindigkeit $v_0$ ab.

    b) Wann erreicht der Ball den höchsten Punkt?

    c) Berechne die maximale Höhe (Fläche im $v$-$t$-Diagramm!).

    d) Erkläre, warum die Geschwindigkeit negativ wird.

    \item \textbf{Vom $v$-$t$- zum $s$-$t$-Diagramm:} Gegeben ist das $v$-$t$-Diagramm eines Aufzugs.

\begin{center}
\begin{tikzpicture}[scale=0.7]
    \begin{axis}[
        xlabel={$t$ in s},
        ylabel={$v$ in m/s},
        xmin=0, xmax=12,
        ymin=0, ymax=5,
        grid=major,
        width=10cm, height=4cm,
        axis lines=left,
        xtick={0,2,4,6,8,10,12},
    ]
    % Trapezform
    \addplot[very thick, physik, domain=0:2] {2*x};
    \addplot[very thick, physik, domain=2:8] {4};
    \addplot[very thick, physik, domain=8:10] {4 - 2*(x-8)};
    \end{axis}
\end{tikzpicture}
\end{center}

    a) In welchen Zeitabschnitten beschleunigt, bremst und fährt der Aufzug gleichförmig?

    b) Berechne die Gesamtstrecke (Fläche unter dem Graphen).

    c) Zeichne das zugehörige $s$-$t$-Diagramm. Achte auf:
    \begin{itemize}[leftmargin=1.5em]
        \item Parabel bei Beschleunigung/Bremsen
        \item Gerade bei gleichförmiger Bewegung
        \item Korrekten Übergang der Steigungen
    \end{itemize}
\end{enumerate}

\newpage

% ═══════════════════════════════════════════════════════════════
% KAPITEL 10: CHECKLISTE
% ═══════════════════════════════════════════════════════════════

\section{Checkliste zur Klausurvorbereitung}

\begin{center}
\renewcommand{\arraystretch}{1.6}
\begin{tabular}{|p{10cm}|c|c|c|}
\hline
\textbf{Ich kann ...} & \textbf{sicher} & \textbf{unsicher} & \textbf{üben!} \\
\hline
\multicolumn{4}{|l|}{\textbf{1. Gleichförmige Bewegung}} \\
\hline
... die Formel $v = s/t$ anwenden und umstellen & $\square$ & $\square$ & $\square$ \\
\hline
... km/h in m/s umrechnen und umgekehrt & $\square$ & $\square$ & $\square$ \\
\hline
... $s$-$t$- und $v$-$t$-Diagramme interpretieren \textbf{und zeichnen} & $\square$ & $\square$ & $\square$ \\
\hline
... den Weg aus dem $v$-$t$-Diagramm berechnen (Fläche) & $\square$ & $\square$ & $\square$ \\
\hline
\multicolumn{4}{|l|}{\textbf{2. Beschleunigte Bewegung}} \\
\hline
... die Beschleunigung berechnen & $\square$ & $\square$ & $\square$ \\
\hline
... die Formeln $v = at$ und $s = \frac{1}{2}at^2$ anwenden & $\square$ & $\square$ & $\square$ \\
\hline
... Bremsweg und Bremszeit berechnen & $\square$ & $\square$ & $\square$ \\
\hline
... $s$-$t$- und $v$-$t$-Diagramme der beschleunigten Bewegung zeichnen und interpretieren & $\square$ & $\square$ & $\square$ \\
\hline
\multicolumn{4}{|l|}{\textbf{3. Zusammengesetzte Bewegungen}} \\
\hline
... Bewegungsgleichungen mit $v_0$ und $s_0$ aufstellen & $\square$ & $\square$ & $\square$ \\
\hline
... Überholprobleme lösen (Gleichsetzen) & $\square$ & $\square$ & $\square$ \\
\hline
... $s$-$t$- und $v$-$t$-Diagramme komplexerer zusammengesetzter Bewegungen interpretieren und zeichnen & $\square$ & $\square$ & $\square$ \\
\hline
\multicolumn{4}{|l|}{\textbf{4. Freier Fall}} \\
\hline
... Fallzeit und Aufprallgeschwindigkeit berechnen & $\square$ & $\square$ & $\square$ \\
\hline
... die Formeln des senkrechten Wurfs anwenden & $\square$ & $\square$ & $\square$ \\
\hline
\multicolumn{4}{|l|}{\textbf{5. Luftwiderstand (nur qualitativ)}} \\
\hline
... erklären, warum $F_W$ mit der Geschwindigkeit zunimmt & $\square$ & $\square$ & $\square$ \\
\hline
... die Ursache der Grenzgeschwindigkeit erklären ($F_W = F_G$) & $\square$ & $\square$ & $\square$ \\
\hline
... das $v$-$t$-Diagramm eines Fallschirmspringers zeichnen/interpretieren & $\square$ & $\square$ & $\square$ \\
\hline
... das $s$-$t$-Diagramm eines Fallschirmspringers zeichnen/interpretieren & $\square$ & $\square$ & $\square$ \\
\hline
\end{tabular}
\end{center}

\vspace{1em}

\begin{tcolorbox}[colback=hellblau, colframe=physik, boxrule=1pt, title=\textbf{Tipps für die Klausur}]
\begin{enumerate}
    \item \textbf{Einheiten zweckmäßig umrechnen:} Meist in SI-Einheiten (m, s, m/s statt km/h). Situationsbedingt kann es aber sinnvoll sein, in km und h zu rechnen (z.\,B. bei langen Strecken).

    \item \textbf{Skizze anfertigen:} Bei Bewegungsaufgaben hilft eine Skizze mit Zeitachse.

    \item \textbf{Nachvollziehbarer Rechenweg:}
    \begin{itemize}[leftmargin=1em]
        \item \textbf{Formel} aufschreiben (ggf. umstellen -- was nicht im Tafelwerk steht, muss hergeleitet werden!)
        \item \textbf{Einsetzen} mit Einheiten
        \item \textbf{Ergebnis} mit Einheit
    \end{itemize}

    \item \textbf{Ergebnis prüfen:} Ist das Ergebnis realistisch? Stimmt die Einheit?

    \item \textbf{Diagramme:} Achsenbeschriftung mit Größe und Einheit nicht vergessen!
\end{enumerate}
\end{tcolorbox}

\vspace{1em}

\begin{center}
\textbf{Viel Erfolg bei der Klausur!}
\end{center}

\end{document}
